\section{Conclusion}
In this paper, we introduced ViSTR-Bench, a visual spatial-temporal reasoning benchmark designed to evaluate whether MLLMs can reason from continuous visual cues in dynamic scenes. 
Unlike prior benchmarks that mainly focus on static spatial attributes, low-level temporal perception, or quantitative prediction, ViSTR-Bench emphasizes temporal evidence, reasoning-oriented task design, and qualitative evaluation. It covers four complementary dimensions, with 12 subtasks and 1,093 high-quality video question-answer pairs across tabletop, indoor, and outdoor scenarios. Through comprehensive evaluations of proprietary, open-source, and specialized spatial MLLMs, we find that current models remain substantially below human performance. We hope ViSTR-Bench can serve as a diagnostic testbed for future research on physically grounded multimodal intelligence, embodied AI, and dynamic world modeling.
